\section{Introduction}
\label{sec:intro}

Diffusion and flow-matching transformers~\cite{peebles2023dit,esser2024scaling} generate high-resolution images
and videos~\cite{blackforestlabs2024flux,kong2024hunyuanvideo,wan2025wan}, but sampling is slow: each of tens of steps
runs dozens of transformer blocks over thousands of tokens. Training-free caching exploits the similarity of
consecutive steps~\cite{ma2024deepcache,selvaraju2024fora,liu2025teacache}: on some steps the block stack is skipped
and its output is approximated from the computed steps. A cache thus makes two decisions: \emph{when} to compute, and
\emph{what} to write on the steps it skips.

Most work has focused on the first decision. Adaptive gates accumulate a cheap proxy of the change of the network
output and recompute once it exceeds a threshold~\cite{liu2025teacache,ma2025magcache,bu2026dicache};
\seacache~\cite{chung2026seacache} measures this change after a spectral-evolution-aware (\sea) filter that suppresses
the noise component of the features. For the second decision, most caches reuse the last computed
output~\cite{ma2024deepcache,selvaraju2024fora,liu2025teacache,chung2026seacache}, while forecasting methods extrapolate
it from several computed steps with Taylor expansions~\cite{liu2025taylorseer}, Hermite or Chebyshev
polynomials~\cite{feng2026hicache,han2026spectrum}, multistep formulas~\cite{zheng2026foca}, or Kalman
filters~\cite{rajabi2026speckcast}.

Reuse and forecasting share a structural limitation: they are \emph{open loop}. The filler predicts the skipped steps
from the past alone, and the sampler integrates its errors into the latent. These errors grow with the distance to
the last computed step: on six reference trajectories of our development set, the relative error of the velocity
obtained from a Kalman forecast of the residual grows from 4--5\% at the first skipped step to over 20\% at the
seventh (supplement). When the next step is computed, its block-stack output reveals how wrong the forecast was, but
existing caches use this observation only to predict the \emph{next} gap; the latent, already advanced along erroneous
velocities, is never revisited. Forecasting reduces this drift but does not remove it: as gaps grow, its advantage
over reuse vanishes (\cref{sec:main-results}).
% Source of the 4-5% / >20% numbers: research_v2/techB/diag_summary.txt (delta-0.6 schedule, resid_pe column,
% gaps 25->33 and 33->41, open loop along the 50-step trajectories of COCO-64 prompts 0-5).

We close the loop without evaluating the network again. Our key observation is that a skipped gap is a
\emph{smoothing} problem. With a Kalman filter as the predictor of the block-stack residual, the innovation at the next
computed step~$b$, \ie, the observed minus the forecast residual, and the covariances accumulated over the gap
determine, through the Rauch--Tung--Striebel (RTS) smoother~\cite{rauch1965maximum}, the minimum-mean-square-error
residuals of all skipped steps under the filter's model. Because only the cheap output layers depend on the residual,
we re-integrate the gap through them and continue sampling from the corrected latent. We call this procedure \emph{predictor-corrector gap smoothing} (\pcgs,
\cref{fig:method}). Like the corrector of UniPC~\cite{zhao2023unipc}, it reuses a network evaluation that the sampler
performs anyway; unlike it, it corrects an entire gap of steps whose network outputs were never computed.

Combined with the unchanged \sea gate of \seacache and a damped-trend Kalman forecaster of the block-stack residual,
\pcgs yields a training-free cache that changes neither the network, nor the sampler, nor the number of transformer
evaluations. Our contributions are:
\begin{itemize}[leftmargin=*,itemsep=1pt,topsep=2pt]
  \item We identify the open-loop nature of caching fillers: errors made on skipped steps are integrated into the
  latent and never corrected, and forecasting loses its advantage over reuse as gaps grow.
  \item We propose \pcgs, which casts each skipped gap as a smoothing problem: once the next step is computed, the RTS
  smoother of the forecaster's own state-space model gives the residuals of the gap, which is re-integrated with the
  output layers only. We give the smoother gain in closed form for the damped-trend model under the irregular
  observations of an adaptive gate. \pcgs adds neither transformer evaluations nor hyperparameters.
  \item On \flux and held-out DrawBench prompts~\cite{saharia2022imagen}, our cache reaches 27.72\,dB PSNR with
  respect to the 50-step sampler at 18.1 of 50 evaluations, against 26.29\,dB published for \seacache at 18.4, and
  22.53 against 21.33\,dB at 13. Against our re-run of \seacache in the same setup, the paired gains are 2.39\,dB
  [2.03, 2.75] and 1.27\,dB [1.02, 1.52], of which \pcgs contributes 0.81 and 0.61\,dB. \pcgs improves on the
  predictor at every budget from 9 to 27 evaluations and also improves reuse and linear extrapolation; its gain comes
  from re-integrating the gap rather than from the exact form of the smoother gains. On HunyuanVideo it improves on
  \seacache in the same setup by 3.17 and 1.92\,dB at 24 and 16 evaluations; on Wan2.1 it exceeds the published
  results by 5.07 and 2.07\,dB at 24 and 17 evaluations.
\end{itemize}
Code will be released.
